Генерация юмора является удобным способом изучения ограничений современных методов постобучения языковых моделей. В таких областях, как математика или программирование, сгенерированный ответ часто можно проверить внешней процедурой: решение либо совпадает с правильным ответом, либо программа проходит тесты или не проходит их. Юмор устроен иначе. Шутка может быть грамматически правильной, связной и соответствующей теме, но при этом не производить нужного эффекта. При этом проблема не сводится только к субъективности юмора. Один и тот же запрос может допускать множество разных удачных шуток, тогда как наблюдаемый набор данных обычно содержит лишь небольшую их часть. Поэтому юмор является не только сложной задачей генерации, но и полезным примером для изучения рассуждения и обучения с подкреплением в неверифицируемых доменах \citep{Loakman2025,lemmens2026computational}.

Эта сложность стала особенно важной после появления больших языковых моделей. Ранние системы вычислительного юмора были узкими, но интерпретируемыми: они опирались на шаблоны, лексические ресурсы или конкретные механизмы каламбуров \citep{ritchie2001current,amin2020survey}. Современные языковые модели обладают гораздо более широким покрытием, однако недавние работы показывают, что беглость текста не равна сильной юмористической компетентности. Крупные оценки по-прежнему помещают модельные системы ниже сильных человеческих результатов, а исследования каламбуров и понимания шуток показывают, что модели часто надежнее воспроизводят поверхностную форму юмора, чем его внутренний механизм \citep{Zhang2024HumorAI,xu2024good,cocchieri2025you,loakman2025comparing}. Иными словами, юмор выявляет знакомую слабость современных моделей: они хорошо продолжают вероятный текст, но хуже выполняют открытый поиск по множеству правдоподобных и неожиданных продолжений.

Один из ответов на эту слабость состоит в явном моделировании процесса рассуждения. Современные системы генерации юмора все чаще используют поэтапное построение подсказок, творческий поиск или декомпозицию на основе теории, а не одношаговое завершение текста. Эти подходы мотивированы простым наблюдением: хорошие шутки часто возникают не через выбор первого вероятного продолжения, а через исследование ассоциаций, переосмыслений, контрастов и форм панчлайна \citep{chen2023prompt,zhong2024let,Tikhonov2024,Kim2025,ZhangJ2025}. Результаты выглядят многообещающе, но имеют ограничение. В большинстве таких систем структура рассуждения задается вручную: исследователь определяет число этапов, тип промежуточных представлений и порядок действий модели.

Параллельное направление в постобучении предлагает другую возможность. В верифицируемых доменах обучение с подкреплением недавно использовалось для извлечения развернутого поведения рассуждения напрямую, а не только через ручное построение подсказок или фиксированные процессы \citep{Lightman2023,DeepSeek2025}. Методы без внешнего верификатора, такие как VeriFree и RLPR, дополнительно показывают, что часть этой логики может быть перенесена за пределы математики и программирования, если заменить внешнего верификатора внутренними вероятностными целями \citep{Zhou2025,Yu2025}. Однако для юмора такие формулировки сразу сталкиваются с несоответствием самой природе задачи. Они построены вокруг предположения об одной эталонной ссылке, тогда как юмор естественно является задачей с несколькими эталонными ответами: для запроса вроде ``write a joke about a frog'' существует много приемлемых ответов, и любой набор данных фиксирует только некоторые из них.

Данная работа мотивирована именно этим несоответствием. Ее центральное утверждение состоит в том, что генерацию юмора следует формулировать не как точную имитацию наблюдаемой шутки, а как распределение вероятностной массы по множеству приемлемых шуток при заданном запросе. Этот сдвиг важен и теоретически, и практически. Теоретически он меняет способ определения обучения без внешнего верификатора с подкреплением в творческом домене. Практически он меняет требования к данным: вместо изолированных пар запрос--ответ обучение требует наборов эталонных ответов, связанных с одним запросом и содержащих несколько правдоподобных шуток.

Работа поэтому состоит из двух связанных частей. Первая часть является методологической и теоретической: в ней формулируется расширение обучения с подкреплением без внешнего верификатора на случай нескольких эталонных ответов для генерации юмора и анализируются последствия неполноты наблюдаемых эталонных наборов. Вторая часть является реализационной: она строит процесс данных, необходимый для практического применения такой цели. В текущем репозитории этот процесс уже включает сбор корпуса, генерацию плотных векторных представлений, извлечение ключевых слов, построение запросы, поиск эталонных ответов, генерацию кандидатов и автоматическую оценку. Отсутствующим этапом остается сам цикл обучения. Соответственно, данную работу следует читать как предварительное исследование в точном смысле этого термина: она задает формальную цель и программную инфраструктуру, необходимую для последующего обучающего эксперимента.

Источником данных является объединенный корпус примерно из 664k шуток, собранный из коллекций Short Jokes и rJokes. Полный корпус используется как пул поиска, тогда как его ключевых слов достаточно для построения обучающего набора данных с несколькими эталонными ответами. Целевая экспериментальная постановка является узким и контролируемым: модели Qwen3-1.7B и Qwen3-4B в режимах instruct и thinking будут сравниваться с вариантом thinking, обученным с помощью MRVF; обучение планируется на NVIDIA H200 GPU с использованием модифицированного \texttt{trl} и \texttt{vLLM} для генерации траекторий и вывода базовых моделей.
